\documentclass[11pt,a4paper]{article}
\usepackage[utf8]{inputenc}
\usepackage[T1]{fontenc}
\usepackage[margin=1in]{geometry}
\usepackage{hyperref}
\usepackage{booktabs}
\usepackage{longtable}
\usepackage{array}
\usepackage{enumitem}
\usepackage{xcolor}
\usepackage{listings}
\usepackage{amsmath}
\usepackage{amssymb}
\usepackage{siunitx}
\usepackage{caption}

\hypersetup{
    colorlinks=true,
    linkcolor=blue,
    urlcolor=blue,
}

\lstset{
    basicstyle=\ttfamily\small,
    breaklines=true,
    frame=single,
    backgroundcolor=\color{gray!5},
    columns=fullflexible,
}

\sisetup{
    group-separator = {,},
    group-minimum-digits = 4,
}

\title{%
    Dataset Statistics Report\\
    \large LaRA, CLERK, CUAD, and GreekBarBench\\
    \large Canonical Long-Context RAG Evaluation Suite
}
\author{Canonical RAG Dataset Suite}
\date{\today}

\begin{document}
\maketitle

\section*{Executive Summary}

This report presents descriptive statistics for four canonical datasets
used in a broader evaluation suite for long-context Retrieval-Augmented
Generation (RAG) systems. The datasets span a diverse range of
document lengths, query characteristics, and retrieval tasks, providing
a heterogeneous basis for comparing RAG architectures on realistic
long-context workloads.

The four datasets in scope are:

\begin{itemize}[leftmargin=*]
    \item \textbf{LaRA} � long-document retrieval benchmark.
    \item \textbf{CLERK} � legal case-law retrieval (CLERC).
    \item \textbf{CUAD} � contract clause extraction.
    \item \textbf{GreekBarBench} � Greek legal bar-examination benchmark.
\end{itemize}

All four datasets are drawn from the same canonicalization pipeline,
which produces a uniform envelope:

\begin{lstlisting}
documents.jsonl   # canonical document records
examples.jsonl    # canonical query + ground truth records
validation.json   # per-dataset validation report
manifest.json     # identity and integrity metadata
\end{lstlisting}

The canonical envelope makes all four datasets directly consumable by
any pipeline that adheres to the shared retrieval contract. All
numbers in this report were computed from the canonical files as they
exist on disk.

\section{Context and Purpose}

\subsection{Broader Evaluation Program}

The four datasets in this report are part of a larger, seven-dataset
canonical suite assembled for evaluating and comparing eleven RAG
architectures on long-context retrieval. The suite spans multiple
domains:

\begin{itemize}[leftmargin=*]
    \item long-context general QA (RULER, LongBench~v2, Loong),
    \item long-document retrieval (LaRA),
    \item legal case law (CLERK),
    \item contractual clause extraction (CUAD),
    \item Greek legal reasoning (GreekBarBench).
\end{itemize}

The four datasets covered in this report were selected because they
exhibit a wide range of corpus characteristics: from very long
documents with short queries (LaRA), to short documents with long
queries (CLERK), to medium-sized contracts with many queries (CUAD),
to structured legal questions (GreekBarBench). This heterogeneity is
intentional; the suite is designed so that no single architecture
benefits disproportionately from a single document-length regime.

\subsection{Design Principles}

\begin{enumerate}[leftmargin=*]
    \item \textbf{Canonical identity.} Every document and every example
    has a stable, deterministic identifier that is preserved across
    all downstream stages.
    \item \textbf{Uniform ground truth.} Every example exposes a
    ground-truth structure with documented cardinality semantics,
    including datasets with legitimately empty ground truth (CUAD
    clause queries that do not appear in a given contract).
    \item \textbf{Query-to-document linkage.} Every query carries an
    explicit reference to its authoritative source documents via
    \texttt{query.doc\_ids}, eliminating ambiguity during retrieval
    evaluation.
    \item \textbf{Document-level preservation.} The canonical
    document text is preserved verbatim. No normalization, no
    chunking, and no embedding occurs at this stage.
\end{enumerate}

\subsection{Intended Use}

The canonical files described here are \emph{pre-retrieval} inputs.
They are intended to be consumed by RAG pipelines that will:

\begin{enumerate}[leftmargin=*]
    \item parse and chunk the documents according to their own
    architecture-specific policy,
    \item generate embeddings with a frozen, revision-pinned encoder
    (in the controlled comparison track),
    \item persist the resulting index,
    \item execute retrieval against the canonical queries,
    \item project ranked items back to canonical document identifiers,
    \item evaluate retrieval quality against the ground truth.
\end{enumerate}

\section{Methodology}

\subsection{Sources}

All statistics in this report were computed directly from the
canonical files under:

\begin{lstlisting}
C:/Users/USER1/Documents/Project 1/processed/canonical/<dataset>/
  documents.jsonl
  examples.jsonl
  validation.json
  manifest.json
\end{lstlisting}

\subsection{Definitions}

\begin{itemize}[leftmargin=*]
    \item \textbf{Document.} One canonical record in
    \texttt{documents.jsonl}. Text is read from
    \texttt{text\_normalized} with fallback to \texttt{text\_raw} or
    \texttt{text}.
    \item \textbf{Example.} One canonical query with its associated
    ground truth, read from \texttt{examples.jsonl}.
    \item \textbf{Token.} Whitespace-delimited token
    (\texttt{len(text.split())}). This is a stable,
    tokenizer-independent count suitable for cross-dataset comparison.
    \item \textbf{Chunk estimate.} A projection of how many chunks a
    document would yield under a fixed chunk size with a fixed
    overlap. The formula is:
    \begin{equation*}
        C = 1 + \left\lceil \frac{T - S + \text{stride} - 1}{\text{stride}} \right\rceil,
        \quad \text{stride} = S - O,
    \end{equation*}
    where $T$ is the token count, $S$ is the chunk size, and $O$ is
    the overlap. Chunking is \emph{not} executed at this stage; the
    numbers are projections only.
    \item \textbf{Ground-truth cardinality.} The number of document
    identifiers listed in \texttt{ground\_truth.relevant\_ids} for one
    example.
\end{itemize}

\subsection{Non-Goals}

This report does \emph{not} include:

\begin{itemize}[leftmargin=*]
    \item tokenizer-specific token counts (e.g.\ BGE-M3 or Qwen
    tokenization),
    \item actual chunked outputs,
    \item embeddings,
    \item retrieval results or evaluation metrics,
    \item any normalization, filtering, or modification of the source
    text.
\end{itemize}

\section{Dataset Profiles}

\subsection{LaRA (Long-context Retrieval Assessment)}

LaRA is a long-document retrieval benchmark designed to stress-test
RAG systems on very long source documents with short, targeted
queries. In LaRA, a single document can easily exceed the
32\,768-token context window of many encoders; the retrieval task is
therefore fundamentally about identifying the correct long document
from a small but very heavy corpus.

\paragraph{Key characteristics.}
\begin{itemize}[leftmargin=*]
    \item \textbf{Corpus size.} 131 documents � the smallest corpus
    in the suite.
    \item \textbf{Document length.} Mean 40{,}463 whitespace tokens;
    maximum 106{,}030. The corpus is dominated by long documents.
    \item \textbf{Query length.} Mean 111 characters � short,
    highly-targeted queries.
    \item \textbf{Ground truth.} 1 relevant document per query
    (cardinality exactly 1).
    \item \textbf{Examples per document.} $\approx 17.8$ queries per
    document on average.
\end{itemize}

\paragraph{What LaRA tests.}
LaRA is the natural stress test for \emph{long-document} retrieval:
can an architecture locate a single needle in a heavily-tokenized
document without relying on chunk-level shortcuts? Because the
ground-truth cardinality is exactly 1, LaRA measures pure
long-document identification.

\subsection{CLERK (CLERC Legal Retrieval)}

CLERK is derived from the CLERC benchmark for legal case-law
retrieval. Unlike LaRA, CLERK pairs short documents with long,
information-rich queries. It is a legal-domain dataset in which the
retrieval task is to identify which passage(s) of a legal corpus are
relevant to a full case description.

\paragraph{Key characteristics.}
\begin{itemize}[leftmargin=*]
    \item \textbf{Corpus size.} 2{,}756 documents � the largest
    corpus in this report.
    \item \textbf{Document length.} Mean 233 whitespace tokens;
    maximum 254. Documents are essentially uniform in length and very
    short.
    \item \textbf{Query length.} Mean 16{,}322 characters � by far
    the longest queries in the suite.
    \item \textbf{Ground truth.} Mean 2.77 relevant documents per
    query; maximum 9. This is the only dataset in this report with
    non-trivial ground-truth cardinality.
    \item \textbf{Examples per document.} $\approx 0.36$ queries per
    document � many documents are not individually referenced.
\end{itemize}

\paragraph{What CLERK tests.}
CLERK places heavy emphasis on \emph{query understanding} and
\emph{multi-document retrieval}. The long queries stress the encoder's
ability to summarize legal facts; the short documents favor
architectures that can perform efficient dense retrieval over a large
pool. Because ground-truth cardinality varies from 1 to 9, CLERK
rewards recall-oriented retrieval policies.

\subsection{CUAD (Contract Understanding Atticus Dataset)}

CUAD is a legal contract analysis dataset. Each document is a
commercial contract; each example is a clause-extraction query
(e.g.\ ``Highlight the parts related to termination for
convenience''). CUAD is characterized by many queries per document,
and by the fact that most clause queries do not apply to most
contracts � the ground truth is legitimately empty for the majority of
query-document pairs.

\paragraph{Key characteristics.}
\begin{itemize}[leftmargin=*]
    \item \textbf{Corpus size.} 510 contracts.
    \item \textbf{Document length.} Mean 7{,}861 whitespace tokens;
    median 5{,}006; maximum 47{,}733.
    \item \textbf{Query length.} Mean 260 characters � short clause
    descriptions.
    \item \textbf{Ground truth.} 20{,}910 examples, of which
    \textbf{14{,}208 (67.9\%) have empty ground truth} and 6{,}702
    have exactly one relevant document. All non-empty examples have
    cardinality 1.
    \item \textbf{Examples per document.} $\approx 41$ clause queries
    per contract.
\end{itemize}

\paragraph{What CUAD tests.}
CUAD stresses \emph{many-shot query evaluation}. A single contract may
be queried dozens of times, once per clause. The dominance of empty
ground truths means any evaluation must include an explicit
\emph{negative} class handling policy: architectures that always
return a top-$k$ result will be penalized if the metric does not
account for empty-relevance queries. CUAD is the primary benchmark
in this suite for measuring false-positive behaviour.

\subsection{GreekBarBench}

GreekBarBench is a Greek-language legal reasoning benchmark built from
bar-examination style questions. Each document is a short legal
excerpt; each query is a bar-exam style question in Greek. It is the
only non-English dataset in this report, providing a multilingual
stress test for the encoder.

\paragraph{Key characteristics.}
\begin{itemize}[leftmargin=*]
    \item \textbf{Corpus size.} 284 documents.
    \item \textbf{Document length.} Mean 10{,}162 whitespace tokens;
    median 9{,}971; maximum 25{,}307. Documents are moderately long
    and relatively uniform.
    \item \textbf{Query length.} Mean 194 characters.
    \item \textbf{Ground truth.} 1 relevant document per query
    (cardinality exactly 1).
    \item \textbf{Examples per document.} Exactly 1 query per
    document.
\end{itemize}

\paragraph{What GreekBarBench tests.}
GreekBarBench is a clean, one-to-one benchmark: one query, one
document, no ambiguity. It measures a pipeline's ability to
correctly identify a single relevant source in a foreign-language
corpus. Its uniform structure makes it the easiest dataset in the
suite for sanity-checking a new architecture, while still being
informative about multilingual encoder behaviour.

\section{Main Statistics}

\begin{table}[h]
\centering
\caption{Main statistics for the four datasets.}
\small
\begin{tabular}{@{}l S[table-format=6.0] S[table-format=6.0] S[table-format=6.0] S[table-format=6.0]@{}}
\toprule
\textbf{Metric} & \textbf{LaRA} & \textbf{CLERK} & \textbf{CUAD} & \textbf{GreekBarBench} \\
\midrule
document count                & 131    & 2756   & 510    & 284    \\
example (GT) count            & 2326   & 1000   & 20910  & 284    \\
total dataset size (bytes)    & 84388337 & 26809056 & 89752275 & 77697117 \\
mean tokens per document      & 40463.66 & 233.49 & 7861.19 & 10161.60 \\
mean chunks per doc @ 512/ov64 & 90.70 & 1.00 & 17.89 & 23.06 \\
\bottomrule
\end{tabular}
\end{table}

\section{File Sizes}

\begin{table}[h]
\centering
\caption{Per-file sizes in bytes.}
\small
\begin{tabular}{@{}l S[table-format=8.0] S[table-format=8.0] S[table-format=8.0] S[table-format=8.0]@{}}
\toprule
\textbf{File} & \textbf{LaRA} & \textbf{CLERK} & \textbf{CUAD} & \textbf{GreekBarBench} \\
\midrule
documents.jsonl & 69390194 & 8577683  & 53378228 & 68949846 \\
examples.jsonl  & 14996106 & 18229557 & 36370354 & 8744487  \\
validation.json & 969      & 702      & 2620     & 1688     \\
manifest.json   & 1068     & 1114     & 1073     & 1096     \\
\bottomrule
\end{tabular}
\end{table}

\section{Document Counts}

\begin{table}[h]
\centering
\caption{Canonical document counts.}
\small
\begin{tabular}{@{}l S[table-format=5.0]@{}}
\toprule
\textbf{Dataset} & \textbf{Documents} \\
\midrule
LaRA          & 131    \\
CLERK         & 2756   \\
CUAD          & 510    \\
GreekBarBench & 284    \\
\bottomrule
\end{tabular}
\end{table}

\section{Example and Ground-Truth Counts}

\begin{table}[h]
\centering
\caption{Example-level counts.}
\small
\begin{tabular}{@{}l S[table-format=6.0] S[table-format=6.0] S[table-format=6.0] S[table-format=6.0]@{}}
\toprule
\textbf{Metric} & \textbf{LaRA} & \textbf{CLERK} & \textbf{CUAD} & \textbf{GreekBarBench} \\
\midrule
example count                     & 2326  & 1000  & 20910 & 284  \\
examples with empty GT            & 0     & 0     & 14208 & 0    \\
examples with GT (non-empty)      & 2326  & 1000  & 6702  & 284  \\
unique relevant documents         & 131   & 2756  & 510   & 284  \\
\bottomrule
\end{tabular}
\end{table}

\section{Ground-Truth Cardinality per Example}

\begin{table}[h]
\centering
\caption{Ground-truth cardinality statistics.}
\small
\begin{tabular}{@{}l S[table-format=2.2] S[table-format=2.2] S[table-format=2.2] S[table-format=2.2]@{}}
\toprule
\textbf{Metric} & \textbf{LaRA} & \textbf{CLERK} & \textbf{CUAD} & \textbf{GreekBarBench} \\
\midrule
mean GT cardinality    & 1.00 & 2.77 & 1.00 & 1.00 \\
median GT cardinality  & 1.00 & 2.00 & 1.00 & 1.00 \\
min GT cardinality     & 1    & 1    & 1    & 1    \\
max GT cardinality     & 1    & 9    & 1    & 1    \\
\bottomrule
\end{tabular}
\end{table}

\section{Token Statistics per Document}

\begin{table}[h]
\centering
\caption{Whitespace-token statistics per document.}
\small
\begin{tabular}{@{}l S[table-format=8.0] S[table-format=8.0] S[table-format=8.0] S[table-format=8.0]@{}}
\toprule
\textbf{Metric} & \textbf{LaRA} & \textbf{CLERK} & \textbf{CUAD} & \textbf{GreekBarBench} \\
\midrule
total whitespace tokens  & 5300739 & 643497 & 4009206 & 2885893  \\
mean tokens per doc      & 40463.66 & 233.49 & 7861.19 & 10161.60 \\
median tokens per doc    & 41179   & 253    & 5006    & 9971    \\
min tokens per doc       & 7695    & 4      & 109     & 2       \\
max tokens per doc       & 106030  & 254    & 47733   & 25307   \\
p95 tokens per doc       & 84447.50 & 253   & 25203.05 & 23879.00 \\
p99 tokens per doc       & 102756.70 & 254  & 41869.66 & 25307.00 \\
\bottomrule
\end{tabular}
\end{table}

\section{Character and Word Statistics per Document}

\begin{table}[h]
\centering
\caption{Character and word statistics per document.}
\small
\begin{tabular}{@{}l S[table-format=9.0] S[table-format=9.0] S[table-format=9.0] S[table-format=9.0]@{}}
\toprule
\textbf{Metric} & \textbf{LaRA} & \textbf{CLERK} & \textbf{CUAD} & \textbf{GreekBarBench} \\
\midrule
total chars             & 33386478 & 3924950 & 25833097 & 18882132 \\
mean chars per doc      & 254858.61 & 1424.15 & 50653.13 & 66486.38 \\
median chars per doc    & 276788  & 1536    & 32152.50 & 67362  \\
min chars per doc       & 49257   & 15      & 645     & 18     \\
max chars per doc       & 624923  & 1793    & 300418  & 164780 \\
mean words per doc      & 40463.66 & 233.49 & 7861.19 & 10161.60 \\
\bottomrule
\end{tabular}
\end{table}

\section{Chunk Estimates per Document}

\begin{table}[h]
\centering
\caption{Projected mean chunks per document, under five chunking configurations.}
\small
\begin{tabular}{@{}l S[table-format=4.2] S[table-format=4.2] S[table-format=4.2] S[table-format=4.2]@{}}
\toprule
\textbf{Configuration} & \textbf{LaRA} & \textbf{CLERK} & \textbf{CUAD} & \textbf{GreekBarBench} \\
\midrule
size 256,  overlap 0    & 158.54 & 1.00 & 31.22 & 40.17 \\
size 512,  overlap 0    & 79.51  & 1.00 & 15.86 & 20.33 \\
size 512,  overlap 64   & 90.70  & 1.00 & 17.89 & 23.06 \\
size 1024, overlap 0    & 39.99  & 1.00 & 8.18  & 10.38 \\
size 1024, overlap 128  & 45.50  & 1.00 & 9.12  & 11.69 \\
\bottomrule
\end{tabular}
\end{table}

\section{Query Statistics}

\begin{table}[h]
\centering
\caption{Query-level statistics.}
\small
\begin{tabular}{@{}l S[table-format=8.2] S[table-format=8.2] S[table-format=8.2] S[table-format=8.2]@{}}
\toprule
\textbf{Metric} & \textbf{LaRA} & \textbf{CLERK} & \textbf{CUAD} & \textbf{GreekBarBench} \\
\midrule
mean query chars                & 111.77   & 16322.28 & 260.00 & 194.16 \\
queries with query.doc\_ids     & 2326     & 1000     & 20910  & 284    \\
total query.doc\_id references  & 2326     & 2770     & 20910  & 284    \\
\bottomrule
\end{tabular}
\end{table}

\section{Aggregate Summary}

\begin{table}[h]
\centering
\caption{Aggregate statistics across the four datasets.}
\begin{tabular}{@{}l S[table-format=8.0]@{}}
\toprule
\textbf{Aggregate metric} & \textbf{Value} \\
\midrule
total documents (all four)                       & 3681       \\
total examples (all four)                        & 24520      \\
total examples with non-empty GT                 & 10312      \\
total examples with empty GT                     & 14208      \\
total dataset size (bytes)                       & 278646785  \\
total whitespace tokens (documents)              & 12839335   \\
\bottomrule
\end{tabular}
\end{table}

\begin{table}[h]
\centering
\caption{Cross-dataset extremes.}
\begin{tabular}{@{}l l@{}}
\toprule
\textbf{Descriptor} & \textbf{Dataset} \\
\midrule
largest corpus                       & CLERK (2756 docs) \\
smallest corpus                      & LaRA (131 docs) \\
most examples                        & CUAD (20910) \\
fewest examples                      & GreekBarBench (284) \\
longest documents (mean)             & LaRA (40463.66 tokens) \\
shortest documents (mean)            & CLERK (233.49 tokens) \\
longest queries (mean characters)    & CLERK (16322.28) \\
shortest queries (mean characters)   & LaRA (111.77) \\
highest mean GT cardinality          & CLERK (2.77) \\
most empty-GT examples               & CUAD (14208) \\
\bottomrule
\end{tabular}
\end{table}

\section{Cross-Dataset Observations}

\subsection{Corpus and Query Length Inversion}

The four datasets display a striking inversion between document and
query length:

\begin{itemize}[leftmargin=*]
    \item \textbf{LaRA} has very long documents (mean 40k tokens) and
    very short queries (mean 112 characters).
    \item \textbf{CLERK} has very short documents (mean 233 tokens)
    and very long queries (mean 16k characters).
    \item \textbf{CUAD} sits in the middle for both.
    \item \textbf{GreekBarBench} has medium-length documents and short
    queries.
\end{itemize}

This inversion is deliberate in the suite design: it ensures that no
single encoder or retrieval policy can dominate all four datasets
through document-length or query-length bias alone.

\subsection{Empty Ground Truth}

CUAD is the only dataset in this report with legitimate empty ground
truths; 67.9\% of CUAD examples have zero relevant documents. Any
retrieval evaluation that includes CUAD must handle the negative class
explicitly. This is a distinct property from the three other datasets,
which all have strictly positive ground truth.

\subsection{Query-to-Document Ratio}

The ratio of queries to documents varies by more than two orders of
magnitude:

\begin{itemize}[leftmargin=*]
    \item LaRA: $\approx 17.8$ queries per document.
    \item GreekBarBench: $1$ query per document (perfectly balanced).
    \item CUAD: $\approx 41$ queries per document.
    \item CLERK: $\approx 0.36$ queries per document.
\end{itemize}

The LaRA and CUAD ratios reward evaluation designs that can
efficiently amortize a single document's indexing cost across many
queries. CLERK's ratio rewards query-side efficiency, and
GreekBarBench is the only strictly balanced dataset in the suite.

\subsection{Cross-Language Coverage}

GreekBarBench is the only non-English dataset in this report. Any
encoder used for the full suite must be multilingual; this is one
reason the controlled comparison track uses BGE-M3 as its shared
encoder.

\section{Limitations}

\begin{itemize}[leftmargin=*]
    \item Token counts are whitespace-based and are \emph{not}
    tokenizer-specific. Real tokenizer counts (e.g.\ BGE-M3 or Qwen
    tokenization) will differ and must be measured separately during
    the embedding stage.
    \item Chunk counts are pure projections. Actual chunking is
    deferred to the architecture-specific preprocessing stage.
    \item Ground-truth cardinality is measured at the
    \emph{document} level. Some datasets carry additional
    span-level, choice-level, or answer-level annotations that are
    not summarized in this report.
    \item The full canonicalization pipeline includes validation
    reports and per-file checksums; this report is a
    \emph{statistical} view, not a substitute for the formal
    validation artifacts.
    \item Only four of the seven canonical datasets are covered here.
    RULER, Loong, and LongBench~v2 are outside the scope of this
    report.
\end{itemize}

\section{Summary}

The four datasets in this report constitute a heterogeneous long-
context RAG evaluation basis:

\begin{itemize}[leftmargin=*]
    \item \textbf{LaRA} measures long-document retrieval with
    single-document ground truth.
    \item \textbf{CLERK} measures legal multi-document retrieval with
    long queries and moderate ground-truth cardinality.
    \item \textbf{CUAD} measures clause-level contract retrieval with
    a dominant negative class.
    \item \textbf{GreekBarBench} measures multilingual single-document
    retrieval with a one-to-one query-document structure.
\end{itemize}

Together they cover a broad range of document lengths
($2$--$106{,}030$ tokens), query lengths
($112$--$16{,}322$ characters), ground-truth cardinalities
($0$--$9$), and languages (English and Greek). This range is
intentional and is what makes the suite suitable for a controlled,
reproducible, architecture-agnostic evaluation of long-context RAG
systems.

\end{document}